\documentclass[11pt]{article}
\usepackage{amsmath,amssymb,amsthm}
\usepackage{booktabs}
\usepackage[margin=1in]{geometry}

\newtheorem{theorem}{Theorem}
\newtheorem{lemma}[theorem]{Lemma}

\title{Problem 662: Fibonacci Paths}
\author{Project Euler Solutions}
\date{}

\begin{document}
\maketitle

\section{Problem Statement}

From $(0,0)$, Alice steps to $(a+x, b+y)$ with $x,y \ge 0$ and $\sqrt{x^2+y^2}$ a Fibonacci number. Let $F(W,H)$ count paths to $(W,H)$. Given $F(3,4)=278$ and $F(10,10)=215846462$, find $F(10000,10000) \bmod 10^9+7$.

\section{Mathematical Foundation}

\begin{lemma}[Step Enumeration]
For Fibonacci number $F_k$, the valid steps are all $(x,y) \in \mathbb{Z}_{\ge 0}^2$ with $x^2+y^2 = F_k^2$. This always includes $(0,F_k)$ and $(F_k,0)$, plus Pythagorean decompositions.
\end{lemma}

\begin{proof}
The axis-aligned pairs satisfy the norm condition trivially. Additional solutions correspond to representations of $F_k^2$ as a sum of two positive squares.
\end{proof}

\begin{theorem}[Path Counting Recurrence]
With $\mathcal{S} = \bigcup_k \mathcal{S}_k$ the full step set:
\[
F(x,y) = \sum_{(dx,dy) \in \mathcal{S}} F(x-dx,\, y-dy), \qquad F(0,0) = 1,\quad F(x,y)=0 \text{ for } x<0 \text{ or } y<0.
\]
\end{theorem}

\begin{proof}
Every path to $(x,y)$ ends with a final step $(dx,dy) \in \mathcal{S}$. The number of such paths equals $F(x-dx,y-dy)$. Since $dx,dy \ge 0$, the quantity $x+y$ strictly decreases along path prefixes, establishing a valid topological order. Summing over all possible last steps gives the recurrence.
\end{proof}

\begin{lemma}[Finiteness]
The number of relevant Fibonacci numbers is $O(\log_\varphi N)$ where $N = \max(W,H)$, so $|\mathcal{S}| = O(\log N)$.
\end{lemma}

\begin{proof}
A valid step requires $F_k \le \sqrt{W^2+H^2}$. Since $F_k \sim \varphi^k/\sqrt{5}$, only $O(\log_\varphi N)$ Fibonacci numbers are relevant.
\end{proof}

\section{Algorithm}

Standard 2D dynamic programming over $(x,y)$ in row-major order, summing contributions from all steps in $\mathcal{S}$, modulo $10^9+7$.

\section{Complexity Analysis}

\begin{itemize}
\item \textbf{Time:} $O(WH|\mathcal{S}|) = O(WH\log N)$. For $W=H=10000$: $\sim 4 \times 10^9$ operations.
\item \textbf{Space:} $O(D \cdot H)$ with rolling array, where $D = \max_{(dx,dy)\in\mathcal{S}} dx$.
\end{itemize}

\section{Answer}

\[
\boxed{860873428}
\]

\end{document}
